% Generated draft: requires mathematical and source review.
\documentclass[11pt]{article}
\usepackage[margin=25mm]{geometry}
\usepackage{fontspec}
\setmainfont{Latin Modern Roman}
\usepackage{amsmath,amssymb,mathtools,unicode-math}
\setmathfont{Latin Modern Math}
\usepackage{enumitem,needspace,xurl,hyperref}
\setlength{\parindent}{0pt}
\setlength{\parskip}{5pt}
\setlength{\emergencystretch}{3em}
\newcommand{\N}{\mathbb{N}}
\newcommand{\Z}{\mathbb{Z}}
\newcommand{\Q}{\mathbb{Q}}
\newcommand{\R}{\mathbb{R}}
\newcommand{\C}{\mathbb{C}}
\newcommand{\F}{\mathbb{F}}
\newcommand{\A}{\mathbb{A}}
\newcommand{\PP}{\mathbb{P}}
\newcommand{\E}{\mathbb{E}}
\newcommand{\eps}{\varepsilon}
\newcommand{\dd}{\,\mathrm{d}}
\begin{document}
\section*{UnsolvedMath ID: 800131}
\textbf{Source problem code:} catalogue entry
\subsection{Definitions and mathematical statement}
Let \(X=B\overline{\Gamma}^{\omega}_{1}\), using the notation of the supplied source, and let \(S^{1}=\{z\in\C:|z|=1\}\). A loop in \(X\) is a continuous map \(\ell:S^{1}\to X\). It is homotopically nontrivial if it is not homotopic to a constant map. Two loops \(\ell,\tau:S^{1}\to X\) are freely homotopic if there exists a continuous map \(H:S^{1}\times[0,1]\to X\) such that \(H(z,0)=\ell(z)\) and \(H(z,1)=\tau(z)\) for every \(z\in S^{1}\); no basepoint is required to remain fixed.

Let \(\mathcal T\) denote the collection of loops in \(X\) that are transverse to the \(\Gamma^{\omega}_{1}\)-structure intended by the source. This is a source-dependent predicate: the supplied record does not define the structure, its realization on \(X\), or the regularity and transversality conditions imposed on these loops.

\textbf{Question 72 (Nariman and Yazdi, 2025).} Does there exist a continuous loop \(\ell:S^{1}\to X\) such that
\[
\ell\text{ is not null-homotopic}
\quad\text{and}\quad
\forall\tau\in\mathcal T,\quad \ell\text{ is not freely homotopic to }\tau?
\]

The source precedes this question with the assertion that a loop in \(B\Gamma_H\) which is freely homotopic to a loop transverse to the \(\Gamma^{\omega}_{1}\)-structure represents a nontrivial element of infinite order in
\[
\pi_1\bigl(B\overline{\Gamma}^{\omega}_{1}\bigr)=\Gamma_H.
\]
This assertion is retained as supplied background, without independent verification. The record does not define \(\Gamma_H\), specify the basepoint for this fundamental group, or explain the relationship between \(B\Gamma_H\), used in the background assertion, and \(B\overline{\Gamma}^{\omega}_{1}\), used in the question. Consequently, the displayed question is a formulation relative to the source's intended spaces and transversality predicate; those geometric identifications remain unresolved.
\subsection{Short English statement}
The question asks whether \(B\overline{\Gamma}^{\omega}_{1}\) contains a non-null-homotopic loop whose free homotopy class has no transverse representative. The supplied background says that classes admitting transverse representatives have infinite order, but it uses a different space notation whose relationship to the question's space is unspecified.
\subsection{Sources}
\begin{itemize}
\item Local UnsolvedMath registry, permanent ID \(800131\), source problem code \texttt{catalogue entry}, title: \emph{Foliations and Diffeomorphism Groups --- Question 72}. The registry attributes the problem to Nariman and Yazdi, with proposed year 2025, and labels it open; this is a supplied status label, not an independently verified assessment.
\item Nariman and Yazdi, \emph{Problem List on Foliations and Diffeomorphism Groups} (2025), Question 72. Supplied source URL: \url{https://arxiv.org/abs/2503.18274}. The provenance records access on 2026-10-16 and extraction by \texttt{arxiv\_latex}; these are supplied provenance details, not a claim of access or verification for this draft.
\item Supplied upstream revision: \texttt{d371ee917ac9647fbb06cd6fc529deec6f3dda82}.
\end{itemize}
\end{document}
